\documentclass[11pt]{article}
\usepackage[margin=1in]{geometry}
\usepackage{amsmath,amssymb,amsthm,mathtools}
\usepackage{mathrsfs}
\usepackage{hyperref}
\usepackage{array}
\usepackage{booktabs}
\usepackage{enumitem}
\usepackage{longtable}
\usepackage{microtype}
\hypersetup{colorlinks=true,linkcolor=blue,citecolor=blue,urlcolor=blue}

\newtheorem{theorem}{Theorem}[section]
\newtheorem{proposition}[theorem]{Proposition}
\newtheorem{lemma}[theorem]{Lemma}
\newtheorem{corollary}[theorem]{Corollary}
\newtheorem{definition}[theorem]{Definition}
\newtheorem{conjecture}[theorem]{Conjecture}
\newtheorem{construction}[theorem]{Construction}
\theoremstyle{remark}
\newtheorem{remark}[theorem]{Remark}
\newtheorem{warning}[theorem]{Warning}

\DeclareMathOperator{\Spec}{Spec}
\DeclareMathOperator{\End}{End}
\DeclareMathOperator{\Tr}{Tr}
\DeclareMathOperator{\tr}{tr}
\DeclareMathOperator{\Hom}{Hom}
\DeclareMathOperator{\Aut}{Aut}
\DeclareMathOperator{\Irr}{Irr}
\DeclareMathOperator{\FPdim}{FPdim}
\DeclareMathOperator{\rk}{rk}
\DeclareMathOperator{\CT}{CT}
\DeclareMathOperator{\diag}{diag}
\DeclareMathOperator{\coker}{coker}
\DeclareMathOperator{\im}{im}
\DeclareMathOperator{\Res}{Res}
\DeclareMathOperator{\ord}{ord}
\DeclareMathOperator{\supp}{supp}
\DeclareMathOperator{\Ind}{Ind}
\DeclareMathOperator{\MW}{MW}

\newcommand{\ZZ}{\mathbb Z}
\newcommand{\QQ}{\mathbb Q}
\newcommand{\RR}{\mathbb R}
\newcommand{\CC}{\mathbb C}
\newcommand{\BB}{\mathbb B}
\newcommand{\FF}{\mathbb F}
\newcommand{\OO}{\mathcal O}
\newcommand{\KK}{\mathcal K}
\newcommand{\eps}{\varepsilon}
\newcommand{\del}{\delta}
\newcommand{\tauz}{\tau}
\newcommand{\wt}{\operatorname{wt}}
\newcommand{\one}{\mathbf 1}

\title{Primitive Support Defects, A-Type Kernels, and Orbifold--CUE--F-Theory Realizations}
\author{}
\date{Working draft, April 2026}

\begin{document}
\maketitle

\begin{abstract}
This paper extracts a common algebraic mechanism from a split-zero support semiring, semilocal adelic boundary extensions, cyclic orbifold fusion lines, CUE derivative moments, Eisenstein--Niemeier lattice transforms, and elliptic-fibration gauge enhancement.  The basic object is the negative of a rank, dimension, Frobenius--Perron dimension, or augmentation character.  For a finite set of primitive supported sectors $I$, the defect object is
\[
        \ZZ^I \xrightarrow{-\Sigma} \ZZ,
        \qquad (a_i)_{i\in I}\longmapsto -\sum_{i\in I}a_i.
\]
Its kernel is the virtual dimension-zero lattice.  If $|I|=r$, this kernel is the root lattice $A_{r-1}$, with the $r=1$ case giving the rootless primitive $A_0=0$ and the $r=3$ case giving the Eisenstein plane $A_2$.  The split-zero globalization semiring supplies the support-theoretic origin of this construction by separating the unsupported zero $\tau$ from the supported arithmetic zero $0_R$.  The semilocal crossed-product boundary supplies analytic $K$-theoretic realizations of the primitive $-1$ charge at finite-prime and archimedean boundary edges.  Cyclic simple-current lines realize the same map as the negative augmentation of $\ZZ[C_r]$, while the full cyclic orbifold double $C_r\times\widehat C_r$, after a Heisenberg cocycle twist, yields the Weyl-operator basis of $M_r(\CC)$ and hence the commutator Lie algebra $\mathfrak{sl}_r$.  CUE matrix-size integration gives the $A_{N-1}$ Weyl denominator, and first-derivative $s$-fold CUE moments carry the exponent $s^2+2s=\dim\mathfrak{sl}_{s+1}$, realized functorially by the block decomposition $\End(\CC^s)\oplus\CC^s\oplus(\CC^s)^\vee\cong\mathfrak{sl}_{s+1}$.  For $r=3$, the augmentation ideal of $\ZZ[C_3]$ is the Eisenstein $A_2$ lattice, whose Epstein zeta is $2^{-s}6\zeta(s)L(s,\chi_{-3})$; twelve such planes globalize to the $12A_2$ Niemeier lattice, adjacent to the Leech lattice and $\ZZ_3$-orbifold moonshine.  The same $A_{r-1}$ root mechanism is the local singularity-to-$SU(r)$ enhancement in F-theory, while Mordell--Weil rank supplies the abelian $U(1)$ rank side.  The paper gives explicit $3\times3$ Weyl matrices for the $r=3$ case, Cartan matrices through $A_3$, and theorem-proof derivations for the finite-prime boundary map, the split-zero zero taxonomy, the fusion diagonalization of the defect character, and the CUE-to-Lie block functor.
\end{abstract}

\tableofcontents

\section{Introduction and main construction}

The recurrent integers
\[
        -1,
        \quad -2,
        \quad -3,
        \quad \ldots
\]
are not best regarded as a numerological tower.  They are the scalar shadows of one additive character.  In every realization considered below, there is a free abelian group generated by primitive irreducible sectors, and there is a canonical dimension-like homomorphism sending each primitive sector to $1$.  The project defect is the negative of that homomorphism.

The same pattern appears in several guises:
\begin{enumerate}[label=(\roman*)]
\item the augmentation character on a finite set or group ring;
\item the rank map on $K_0$ of a commutative boundary algebra;
\item the ordinary dimension character on a representation ring;
\item the Frobenius--Perron dimension character on a fusion ring;
\item the support-cardinality character obtained by linearizing Boolean support in the split-zero semiring;
\item the boundary connecting map in local semilocal crossed-product extensions, once the boundary orientation is fixed.
\end{enumerate}
The stable object is therefore
\[
        \delta=-d,
\]
where $d$ denotes the relevant rank, dimension, or augmentation character.

The first theorem is the common algebraic core.

\begin{theorem}[primitive support defect]
Let $X$ be a finite set.  Let
\[
        D(X)=\ZZ[X]
\]
be the free abelian group on $X$, and let
\[
        \epsilon_X:D(X)\to\ZZ,
        \qquad
        \epsilon_X\left(\sum_{x\in X}a_x[x]\right)=\sum_{x\in X}a_x.
\]
Define
\[
        \delta_X=-\epsilon_X.
\]
Then every primitive sector has defect $-1$:
\[
        \delta_X([x])=-1.
\]
If $|X|=r$, then
\[
        \delta_X\left(\sum_{x\in X}[x]\right)=-r.
\]
Moreover,
\[
        \ker\delta_X=\ker\epsilon_X\cong A_{r-1},
\]
with the convention $A_0=0$ when $r=1$.
\end{theorem}

\begin{proof}
Choose an ordering $X=\{1,\ldots,r\}$ and identify $D(X)$ with $\ZZ^r$ with basis $e_1,\ldots,e_r$.  Then
\[
        \delta_X(a_1,\ldots,a_r)=-(a_1+\cdots+a_r).
\]
The kernel is
\[
        K_r=\{(a_1,\ldots,a_r)\in\ZZ^r: a_1+\cdots+a_r=0\}.
\]
For $r\geq2$, put
\[
        \alpha_i=e_i-e_{i+1},\qquad 1\leq i\leq r-1.
\]
Each $\alpha_i\in K_r$.  Conversely, if $a=(a_1,\ldots,a_r)\in K_r$ and
\[
        b_i=a_1+\cdots+a_i,
\]
then
\[
        a=\sum_{i=1}^{r-1}b_i(e_i-e_{i+1}).
\]
Thus the $\alpha_i$ form a $\ZZ$-basis of $K_r$.  Under the standard Euclidean form,
\[
        \langle\alpha_i,\alpha_j\rangle=
        \begin{cases}
        2,& i=j,\\
        -1,& |i-j|=1,\\
        0,& |i-j|>1.
        \end{cases}
\]
This is the Cartan matrix of type $A_{r-1}$.  For $r=1$, the kernel is zero, denoted $A_0$ in the present indexing.\end{proof}

The primitive $-1$ layer is therefore not a meaningless zero case.  It is the dimension-one primitive itself.  Its cancellation kernel is zero, but the defect map is nonzero:
\[
        \ZZ\xrightarrow{-1}\ZZ.
\]
This is the algebraic form of the rootless $U(1)$ primitive.  The higher root systems are built by taking differences between such primitive $U(1)$ charges.

\section{Split-zero globalization and the taxonomy of zero}

The split-zero construction begins with a commutative ring $R$ and a semiring
\[
        G(R)=R\sqcup\{\tau\}.
\]
The intended point is not merely that a new symbol has been added.  The point is that the ordinary ring zero $0_R$ remains present as a supported arithmetic element, while $\tau$ becomes the unsupported semiring zero.  The operations are
\[
        r\oplus s=r+s,
        \qquad
        r\oplus\tau=r,
        \qquad
        \tau\oplus\tau=\tau,
\]
for $r,s\in R$, and
\[
        rs=rs,
        \qquad
        r\tau=\tau r=\tau,
        \qquad
        \tau^2=\tau.
\]
The secondary split-zero notes identify the equivalent coordinate model
\[
        G(R)\cong D_R:=\{(r,b)\in R\times\BB: b=0\Rightarrow r=0_R\}
        =\{(0_R,0)\}\cup(R\times\{1\}),
\]
with
\[
        \tau=(0_R,0),
        \qquad
        0_R=(0_R,1).
\]
The same source records the universal ring quotient and unique Boolean support character
\[
        p_R:G(R)\to R,
        \qquad
        p_R(\tau)=0_R,
        \quad p_R(r)=r,
\]
\[
        \chi_R:G(R)\to\BB,
        \qquad
        \chi_R(\tau)=0,
        \quad \chi_R(r)=1\quad(r\in R),
\]
with the semimodule category over $G(R)$ classified by support diagrams and free semimodules described by Boolean-cubical support data; see the split-zero notes and the extracted theorem ledger \cite{SplitZeroNotes}.

\begin{definition}[zero statuses]
Inside the coordinate model $D_R\subset R\times\BB$ there are three distinct zero-like statuses:
\begin{align*}
        \tau&=(0_R,0), &&\text{unsupported semiring zero},\\
        0_R&=(0_R,1), &&\text{supported arithmetic zero},\\
        (r,0),\ r\ne0, &&&\text{disallowed unsupported nonzero state}.
\end{align*}
\end{definition}

Thus one must not identify ``zero coefficient'', ``absent support'', and ``trace-removed identity''.  A coefficient equal to $0_R$ is supported zero.  The only unsupported element is $\tau$.  The states $(r,0)$ for $r\ne0$ are not part of the semiring at all.

\begin{proposition}[support and arithmetic zero primes]
Let $R$ be a nonzero commutative ring.  The subset
\[
        \mathfrak s_R=\{\tau\}
\]
is a prime ideal of $G(R)$.  If $R$ is an integral domain, then
\[
        \mathfrak z_R=\{\tau,0_R\}
\]
is also a prime ideal.  Thus in the domain case there is a chain
\[
        \{\tau\}\subset\{\tau,0_R\}
\]
separating the support prime from the supported arithmetic-zero prime.
\end{proposition}

\begin{proof}
The set $\{\tau\}$ is an ideal because $x\tau=\tau x=\tau$ for all $x\in G(R)$ and $\tau\oplus\tau=\tau$.  If $xy=\tau$ and neither factor is $\tau$, then $x,y\in R$, and their product is the ordinary product in $R\subset G(R)$, hence cannot be $\tau$.  Therefore one factor is in $\{\tau\}$, proving primality.

Assume now that $R$ is a domain.  The set $\{\tau,0_R\}$ is closed under addition and absorbs multiplication.  If $xy\in\{\tau,0_R\}$, then either one factor is $\tau$, or both factors lie in $R$ and their ordinary product is $0_R$.  Since $R$ is a domain, one of the factors is $0_R$.  Therefore one factor lies in $\{\tau,0_R\}$.\end{proof}

\begin{construction}[support linearization]
Let $I$ be a finite set.  The free semimodule $G(R)^I$ has a support map
\[
        \supp:G(R)^I\to\mathcal P(I),
        \qquad
        \supp(x)=\{i\in I: \chi_R(x_i)=1\}.
\]
Linearize Boolean support by
\[
        \mathcal L_I:\mathcal P(I)\to\ZZ^I,
        \qquad
        J\mapsto\sum_{j\in J}e_j.
\]
The support defect is
\[
        \delta_I(J)= -|J|.
\]
Equivalently,
\[
        \ZZ^I\xrightarrow{-\Sigma}\ZZ,
        \qquad
        (a_i)\mapsto-\sum_{i\in I}a_i.
\]
\end{construction}

\begin{theorem}[support origin of the $A$-tower]
The linearized split-zero support kernel is
\[
        \ker\left(\ZZ^I\xrightarrow{-\Sigma}\ZZ\right)
        \cong A_{|I|-1}.
\]
Thus the $-1,-2,-3,\ldots$ tower is obtained by repeated primitive supported sectors, while roots are differences between primitive supported sectors.
\end{theorem}

\begin{proof}
This is Theorem 1.1 applied to the finite set $I$.  The split-zero input is the existence of the support map $G(R)^I\to\mathcal P(I)$ and the fact that $\tau$ and $0_R$ are separated by $\chi_R$.\end{proof}

\section{Primitive $U(1)$ charges, Cartan matrices, and recursive roots}

The $U(1)$ primitive appears at every level.  A rank-$r$ primitive sector should be read as $r$ independent $U(1)$-type supported charges before imposing the diagonal relation.  The root lattice is the lattice of differences between these primitive charges.

Let
\[
        T_r=U(1)^r.
\]
Its character lattice is
\[
        X^*(T_r)=\ZZ^r.
\]
The diagonal subtorus
\[
        U(1)_\Delta=\{(z,\ldots,z):z\in U(1)\}
\]
corresponds to the total character.  The quotient torus
\[
        T_r/U(1)_\Delta
\]
is a maximal torus of $PU(r)$, while the determinant-one subtorus
\[
        \{(z_1,\ldots,z_r):z_1\cdots z_r=1\}
\]
is a maximal torus of $SU(r)$.

\begin{theorem}[roots as differences of primitive $U(1)$ charges]
The root lattice of $SU(r)$ is
\[
        A_{r-1}=\{(a_1,\ldots,a_r)\in\ZZ^r: \sum_i a_i=0\}.
\]
Equivalently, it is the virtual dimension-zero lattice of the $r$ primitive $U(1)$ characters.
\end{theorem}

\begin{proof}
Let $\mathfrak h$ be the diagonal Cartan subalgebra of $\mathfrak{sl}_r(\CC)$:
\[
        \mathfrak h=\{\diag(x_1,\ldots,x_r):\sum_i x_i=0\}.
\]
For the elementary matrix $E_{ij}$ and $H=\diag(x_1,\ldots,x_r)$,
\[
        [H,E_{ij}]=(x_i-x_j)E_{ij}.
\]
Thus the roots are $e_i-e_j$, $i\ne j$.  They generate the coordinate-sum-zero lattice, which is $A_{r-1}$.\end{proof}

The first Cartan matrices are as follows:
\[
C(A_0)=(),
\qquad
C(A_1)=\begin{pmatrix}2\end{pmatrix},
\]
\[
C(A_2)=\begin{pmatrix}2&-1\\-1&2\end{pmatrix},
\qquad
C(A_3)=\begin{pmatrix}2&-1&0\\-1&2&-1\\0&-1&2\end{pmatrix}.
\]
The rank-three kernel case is $A_3$, and the standard root-system coincidence
\[
        A_3\cong D_3
\]
identifies the $SU(4)$ and $Spin(6)$ root systems.  Thus the first rank-three cancellation kernel is already a trial point where the $A$-tower touches the $D$-tower.

The standard-model gauge-algebra pattern is visible at the level of this root decomposition:
\[
        A_2\oplus A_1\oplus A_0
        \quad\leadsto\quad
        \mathfrak{su}_3\oplus\mathfrak{su}_2\oplus\mathfrak u_1.
\]
This is only a root datum until a charge lattice and quotient by the finite center are specified.  The actual Standard Model gauge group involves a quotient of $SU(3)\times SU(2)\times U(1)$ by a finite subgroup; the present construction supplies the primitive support and root-lattice skeleton, not the full phenomenological model.

\section{The semilocal boundary extension and the finite-prime edge}

The semilocal adelic construction of Connes--Consani considers, for a finite set of places $S$ containing $\infty$, a visible space obtained from semilocal adeles by retaining the generic stratum and the strata with exactly one zero coordinate.  For $S=\{p,q,\infty\}$, the relevant quotient decomposes as
\[
        Z=Z_\varnothing\sqcup Z_p\sqcup Z_q\sqcup Z_\infty,
\]
and the crossed product
\[
        A=C_0(Z)\rtimes\Gamma,
        \qquad
        \Gamma=\{\pm p^mq^n:m,n\in\ZZ\},
\]
sits in
\[
        0\to A_\varnothing\to A\to A_p\oplus A_q\oplus A_\infty\to0.
\]
The literature computes the Morita types
\[
        A_\varnothing\sim_M \KK\otimes C_0(\RR_+^*),
\]
\[
        A_p\sim_M\KK\otimes C(C_p),
        \qquad
        A_q\sim_M\KK\otimes C(C_q),
\]
\[
        A_\infty\sim_M\KK\otimes C^*(\ZZ/2\ZZ),
\]
and hence obtains $K_0(A)\cong\ZZ^3$, $K_1(A)\cong\ZZ^2$ for $S=\{p,q,\infty\}$; see \cite{ConnesConsaniKnots} and the extracted notes \cite{SplitZeroNotes}.

The componentwise connecting map is finer than the rank computation.  The archimedean component is the familiar fixed-point extension
\[
        C_0(\RR)\rtimes\ZZ/2\ZZ.
\]
The finite-prime component is a one-sided valuation compactification with boundary circle.  We now compute it.

\subsection{The finite local model}

Fix a finite prime $p$ and let
\[
        \ell_p=\log p.
\]
The local valuation coordinate is $m\in\ZZ$, with $m\to+\infty$ corresponding to the zero $p$-adic coordinate.  Let
\[
        Y_p=(\ZZ\cup\{\infty\})\times\RR,
\]
where $\infty$ is the endpoint of $\ZZ$ approached by $m\to+\infty$.  Let $T$ be the generator of $\ZZ$ acting by
\[
        T(m,y)=(m+1,y+\ell_p),
        \qquad
        T(\infty,y)=(\infty,y+\ell_p).
\]
Set
\[
        E_p=C_0(Y_p)\rtimes_T\ZZ,
\]
\[
        I_p=C_0(\ZZ\times\RR)\rtimes_T\ZZ,
\]
\[
        Q_p=C_0(\{\infty\}\times\RR)\rtimes_T\ZZ.
\]
There is a short exact sequence
\[
        0\to I_p\to E_p\to Q_p\to0.
\]

\begin{lemma}[Morita types of the finite edge]
There are Morita equivalences
\[
        I_p\sim_M C_0(\RR),
        \qquad
        Q_p\sim_M C(S^1).
\]
\end{lemma}

\begin{proof}
On the generic part $\ZZ\times\RR$, the action is free and proper.  The invariant coordinate is
\[
        t=y-m\ell_p,
\]
since
\[
        (y+\ell_p)-(m+1)\ell_p=y-m\ell_p.
\]
Thus the orbit space of the generic part is $\RR$.  Green imprimitivity gives $I_p\sim_M C_0(\RR)$.

On the boundary $\{\infty\}\times\RR$, the action is translation $y\mapsto y+\ell_p$, so the orbit space is
\[
        \RR/\ell_p\ZZ\cong S^1.
\]
Again Green imprimitivity gives $Q_p\sim_M C(S^1)$.\end{proof}

The compactified quotient may be described explicitly.  Its generic coordinate is
\[
        t=y-m\ell_p,
\]
and its boundary angle is
\[
        \theta=y\mod \ell_p\ZZ.
\]
If $m\to+\infty$ with $y$ bounded, then $t\to-\infty$ and the limiting angle is $t\mod\ell_p\ZZ$.

\begin{theorem}[finite-prime boundary map]
Under the Morita identifications
\[
        K_0(Q_p)\cong K_0(C(S^1))\cong\ZZ[1],
        \qquad
        K_1(I_p)\cong K_1(C_0(\RR))\cong\ZZ,
\]
the connecting map
\[
        \partial_p:K_0(Q_p)\to K_1(I_p)
\]
satisfies
\[
        \partial_p([1])=-1,
\]
where the generator of $K_1(C_0(\RR))$ is represented by a counterclockwise unit loop after one-point compactification of $\RR$.
\end{theorem}

\begin{proof}
Let
\[
        \phi:\RR\to[0,1]
\]
be continuous with
\[
        \lim_{t\to-\infty}\phi(t)=1,
        \qquad
        \lim_{t\to+\infty}\phi(t)=0.
\]
For example, take $\phi(t)=1/(1+e^t)$.  The function
\[
        h(m,y)=\phi(y-m\ell_p)
\]
on the generic part and
\[
        h(\infty,y)=1
\]
on the boundary is continuous in the quotient compactification because $m\to+\infty$ implies $y-m\ell_p\to-\infty$ for bounded $y$.  It is a self-adjoint lift of the unit projection in $Q_p$.

The exponential formula for the boundary map gives the unitary
\[
        u(t)=\exp(2\pi i\phi(t))
\]
in the unitization of $C_0(\RR)$.  As $t$ increases from $-\infty$ to $+\infty$, $\phi(t)$ decreases from $1$ to $0$.  Hence $u(t)$ winds once clockwise around the unit circle.  With counterclockwise winding as $+1$, this class is $-1$.\end{proof}

\begin{theorem}[componentwise semilocal boundary map]
For the visible semilocal algebra with finite primes $p_1,\ldots,p_n$ and $\infty$, the boundary $K_0$ group has primitive generators
\[
        [1]_{p_1},\ldots,[1]_{p_n},[p_+],[p_-].
\]
With the orientation conventions above,
\[
        \delta_0([1]_{p_i})=-1,
        \qquad
        \delta_0([p_+])=-1,
        \qquad
        \delta_0([p_-])=-1.
\]
Hence
\[
        \delta_0(k_1,\ldots,k_n,n_+,n_-)
        =-\left(\sum_{i=1}^n k_i+n_++n_-\right).
\]
\end{theorem}

\begin{proof}
The finite-prime computation is the preceding theorem, transported to the semilocal edge by the explicit valuation model and Green imprimitivity.  The archimedean edge is the local fixed-point model $C_0(\RR)\rtimes\ZZ/2\ZZ$.  Its quotient at the fixed point is $C^*(\ZZ/2\ZZ)$ with minimal projections
\[
        p_+=\frac{1+U}{2},
        \qquad
        p_-=\frac{1-U}{2}.
\]
Using an even lift $\phi(x)p_\pm$ with $\phi(0)=1$ and $\phi(x)\to0$ as $x\to+\infty$, the determinant of the boundary unitary is again $e^{2\pi i\phi(x)}$, winding once clockwise.  Thus each minimal projection maps to $-1$.  Additivity gives the formula.\end{proof}

\begin{corollary}[semilocal triplet]
For $S=\{p,\infty\}$, the class
\[
        [1]_p+[p_+]+[p_-]
\]
has total boundary defect
\[
        -3.
\]
Its liftable virtual boundary lattice is
\[
        \{(a,b,c)\in\ZZ^3:a+b+c=0\}\cong A_2.
\]
\end{corollary}

This result is the semilocal $K$-theoretic realization of the same augmentation object that appears in $\ZZ[C_3]$.

\section{Critical-line quotient coordinates and the anti-invariant primitive}

Let
\[
        u=s(s-1),
        \qquad
        v=2s-1.
\]
Then
\[
        u(1-s)=u(s),
        \qquad
        v(1-s)=-v(s),
        \qquad
        v^2=1+4u.
\]
The coordinate $u$ descends to the quotient by $s\mapsto1-s$, while $v$ is the anti-invariant lift coordinate.

The principal part space with possible simple poles at $0$ and $1$ has the basis
\[
        \frac1u=\frac1{s-1}-\frac1s,
        \qquad
        \frac vu=\frac1{s-1}+\frac1s.
\]
If
\[
        P(s)=\frac a s+\frac b{s-1},
\]
then
\[
        P(s)=A\frac1u+B\frac vu,
        \qquad
        A=\frac{b-a}{2},
        \quad B=\frac{a+b}{2}.
\]

For the completed shifted zeta family
\[
        \Lambda_t(s)=\pi^{-s/2}\Gamma(s/2)\zeta(s,1+t),
\]
the endpoint residues are
\[
        \Res_{s=0}\Lambda_t=-1-2t,
        \qquad
        \Res_{s=1}\Lambda_t=1.
\]
Therefore
\[
        \operatorname{PP}_{\{0,1\}}(\Lambda_t)
        =(1+t)\frac1u-t\frac vu.
\]
At $t=1$,
\[
        \operatorname{PP}_{\{0,1\}}(\Lambda_1)=2\frac1u-\frac vu.
\]
Thus the primitive anti-invariant coefficient is $-1$.

On the critical line
\[
        s=\frac12+i\gamma,
\]
we have
\[
        u=-\left(\gamma^2+\frac14\right)\in\RR,
        \qquad
        v=2i\gamma.
\]
Hence
\[
        \frac1u\in\RR,
        \qquad
        \frac vu\in i\RR.
\]
The critical line is therefore the locus on which the quotient term and the anti-invariant lift term are orthogonal over $\CC$: one is real and one is purely imaginary.  The branch center $s=1/2$ is the point where $v=0$ and the anti-invariant coordinate vanishes.

The logarithmic differential is
\[
        d\log u=\frac vu\,ds.
\]
Thus
\[
        \Res_0(d\log u)=1,
        \qquad
        \Res_1(d\log u)=1,
        \qquad
        \Res_\infty(d\log u)=-2.
\]
The total $-2$ is the doubled primitive anti-invariant $-1$ carried by the two finite endpoints.

\section{Fusion rings, negative dimension, and diagonalization by the $S$-matrix}

The phrase ``negative dimension'' should be understood precisely.  In a fusion ring, the dimension character is a ring homomorphism to $\RR$.  In a pointed line, it is the augmentation.  The defect is its negative.

Let $C_r=\langle g:g^r=1\rangle$ and let
\[
        R_r=\ZZ[C_r].
\]
The augmentation is
\[
        \epsilon\left(\sum_{a=0}^{r-1}n_ag^a\right)=\sum_{a=0}^{r-1}n_a.
\]
This is also the Frobenius--Perron dimension on the pointed fusion ring generated by a simple current $g$.

\begin{proposition}[cyclic line]
The defect
\[
        \delta=-\epsilon
\]
satisfies
\[
        \delta(1+g+\cdots+g^{r-1})=-r,
\]
and
\[
        \ker\epsilon\cong A_{r-1}.
\]
\end{proposition}

\begin{proof}
The first identity is immediate.  The second is Theorem 1.1 applied to the basis $1,g,\ldots,g^{r-1}$.\end{proof}

Let
\[
        \omega=e^{2\pi i/r}.
\]
Over $\CC$, the Fourier matrix
\[
        S_{ab}=r^{-1/2}\omega^{ab},
        \qquad 0\le a,b<r,
\]
diagonalizes the fusion matrices.  Multiplication by $g$ has eigenvalue $\omega^{-b}$ or $\omega^b$, depending on the left/right convention.  The trivial Fourier character is $b=0$.

\begin{theorem}[defect as negative trivial Fourier row]
In the cyclic pointed fusion algebra, the augmentation $\epsilon$ is the projection onto the trivial Fourier character multiplied by the dimension of the regular basis.  Equivalently, the defect $\delta=-\epsilon$ is the negative trivial $S$-row functional.
\end{theorem}

\begin{proof}
Let
\[
        x=\sum_{a=0}^{r-1}n_ag^a.
\]
The Fourier coefficient at $b=0$ is
\[
        \widehat x(0)=\sum_{a=0}^{r-1}n_a.
\]
But this is exactly $\epsilon(x)$.  Hence $\delta(x)=-\widehat x(0)$.  The nontrivial Fourier coefficients span $\ker\epsilon\otimes\CC$.\end{proof}

The regular element
\[
        \rho_r=1+g+\cdots+g^{r-1}
\]
acts diagonally with eigenvalue $r$ on the trivial character and $0$ on all nontrivial characters.  Thus the negative dimension defect is exactly the negative of the trivial diagonal eigenvalue of the regular element.

For a general modular tensor category, the Frobenius--Perron dimension is similarly detected by the positive Perron--Frobenius eigenvector of fusion matrices, and in unitary modular cases by the $S$-matrix row:
\[
        \operatorname{FPdim}(X_i)=\frac{S_{0i}}{S_{00}}.
\]
The project defect is the negative of this dimension character.  The virtual dimension-zero lattice is its kernel.

\section{Orbifold irreducibles, discrete torsion, and the Weyl-operator double}

Dijkgraaf--Vafa--Verlinde--Verlinde analyze orbifolds of rational conformal field theories by finite groups.  For holomorphic orbifolds, the untwisted fusion algebra is the representation algebra of the finite group, and the full twisted sector structure is governed by conjugacy classes and stabilizer representations, up to phases/discrete torsion \cite{DVVV}.  The older string-orbifold construction of Dixon--Harvey--Vafa--Witten shows geometrically why twisted sectors must be included: strings may close only up to a group action \cite{DixonHarveyVafaWitten}.

For a finite group $G$, the standard full holomorphic orbifold label set is
\[
        \operatorname{OrbIrr}(G)=\{(C,\rho): C\in\operatorname{Cl}(G),\ \rho\in\operatorname{Irr}(Z_G(g_C))\},
\]
where $g_C$ is a representative of the conjugacy class $C$.

If $G$ is abelian, then all conjugacy classes are singletons and all centralizers are $G$, so
\[
        \operatorname{OrbIrr}(G)=G\times\widehat G.
\]
For $G=C_r$, this has $r^2$ labels.

This fusion label set is not itself a Lie algebra.  The Lie algebra appears after a discrete-torsion/Heisenberg-cocycle realization.

Let $V=\CC^r$ with basis $\{|m\rangle:m\in\ZZ/r\ZZ\}$, and set
\[
        X|m\rangle=|m+1\rangle,
        \qquad
        Z|m\rangle=\omega^m|m\rangle.
\]
Then
\[
        ZX=\omega XZ.
\]
For $(a,b)\in\ZZ_r^2$, define
\[
        W_{a,b}=X^aZ^b.
\]

\begin{theorem}[twisted group algebra and $\mathfrak{sl}_r$]
The operators $W_{a,b}$ form a basis of $M_r(\CC)$.  They obey
\[
        W_{a,b}W_{c,d}=\omega^{bc}W_{a+c,b+d},
\]
and
\[
        [W_{a,b},W_{c,d}]
        =\left(\omega^{bc}-\omega^{ad}\right)W_{a+c,b+d}.
\]
The nonidentity operators
\[
        W_{a,b},\qquad (a,b)\ne(0,0),
\]
span $\mathfrak{sl}_r(\CC)$.
\end{theorem}

\begin{proof}
The commutation relation $ZX=\omega XZ$ gives
\[
        Z^bX^c=\omega^{bc}X^cZ^b,
\]
and hence the multiplication formula.  Reversing the two factors gives
\[
        W_{c,d}W_{a,b}=\omega^{ad}W_{a+c,b+d},
\]
so the commutator formula follows.

The $r^2$ operators $X^aZ^b$ are linearly independent because $a$ controls the cyclic shift and $b$ controls the independent Fourier phase on that shift.  Since $\dim M_r(\CC)=r^2$, they form a basis.  The trace of $W_{a,b}$ is zero unless $(a,b)=(0,0)$: if $a\ne0$ it has no diagonal entries, and if $a=0,b\ne0$ then the trace is $\sum_m\omega^{bm}=0$.  Therefore the nonidentity $W_{a,b}$ span the traceless matrices.\end{proof}

The cocycle
\[
        \alpha((a,b),(c,d))=\omega^{bc}
\]
is the discrete-torsion phase producing the twisted group algebra
\[
        \CC_\alpha[C_r\times\widehat C_r]\cong M_r(\CC).
\]
For $r=2$ this cocycle is a sign; for $r=3$ it is genuinely Eisenstein, with phases $1,\omega,\omega^2$.

\section{The explicit $r=3$ matrices}

Let
\[
        \omega=e^{2\pi i/3}.
\]
Use the ordered basis $|0\rangle,|1\rangle,|2\rangle$.  Then
\[
X=\begin{pmatrix}
0&0&1\\
1&0&0\\
0&1&0
\end{pmatrix},
\qquad
Z=\begin{pmatrix}
1&0&0\\
0&\omega&0\\
0&0&\omega^2
\end{pmatrix},
\]
with
\[
        ZX=\omega XZ.
\]
The identity sector is
\[
        W_{0,0}=I_3.
\]
The eight nonidentity Weyl operators are:
\[
W_{1,0}=X=\begin{pmatrix}0&0&1\\1&0&0\\0&1&0\end{pmatrix},
\qquad
W_{2,0}=X^2=\begin{pmatrix}0&1&0\\0&0&1\\1&0&0\end{pmatrix},
\]
\[
W_{0,1}=Z=\begin{pmatrix}1&0&0\\0&\omega&0\\0&0&\omega^2\end{pmatrix},
\qquad
W_{0,2}=Z^2=\begin{pmatrix}1&0&0\\0&\omega^2&0\\0&0&\omega\end{pmatrix},
\]
\[
W_{1,1}=XZ=\begin{pmatrix}0&0&\omega^2\\1&0&0\\0&\omega&0\end{pmatrix},
\qquad
W_{1,2}=XZ^2=\begin{pmatrix}0&0&\omega\\1&0&0\\0&\omega^2&0\end{pmatrix},
\]
\[
W_{2,1}=X^2Z=\begin{pmatrix}0&\omega&0\\0&0&\omega^2\\1&0&0\end{pmatrix},
\qquad
W_{2,2}=X^2Z^2=\begin{pmatrix}0&\omega^2&0\\0&0&\omega\\1&0&0\end{pmatrix}.
\]
They form a basis of $\mathfrak{sl}_3(\CC)$ after removing $I_3$ from the full Weyl basis.

The Cartan subalgebra is the diagonal traceless part, spanned for example by
\[
        H_1=\diag(1,-1,0),
        \qquad
        H_2=\diag(0,1,-1).
\]
The roots are
\[
        e_i-e_j,
        \qquad i\ne j.
\]
The simple-root Cartan matrix is
\[
        C(A_2)=\begin{pmatrix}2&-1\\-1&2\end{pmatrix}.
\]

\begin{remark}[where $\tau$ is not]
Inside $G(M_3(\CC))$, the supported zero matrix $0_{M_3}$ and the unsupported zero $\tau$ are different.  A zero coefficient in a Weyl expansion
\[
        A=\sum_{a,b}c_{a,b}W_{a,b}
\]
means $c_{a,b}=0_\CC$ on the supported sheet.  It does not mean that the sector $W_{a,b}$ is globally absent.  Absence is the support condition represented by $\tau$ in the split-zero extension.  Likewise, removing $W_{0,0}=I$ to form $\mathfrak{sl}_3$ is trace removal, not unsupportedness.
\end{remark}

\section{The $C_3$ augmentation ideal, Eisenstein integers, and the $A_2$ Epstein zeta}

For $C_3=\langle g:g^3=1\rangle$, the group ring
\[
        \ZZ[C_3]=\ZZ\{1,g,g^2\}
\]
has augmentation
\[
        \epsilon(a+bg+cg^2)=a+b+c.
\]
The defect is $-\epsilon$ and
\[
        -\epsilon(1+g+g^2)=-3.
\]
The augmentation ideal
\[
        I_{C_3}=\ker\epsilon
\]
has basis
\[
        \alpha_1=1-g,
        \qquad
        \alpha_2=g-g^2.
\]
The Gram matrix in the coefficient inner product is
\[
        \begin{pmatrix}2&-1\\-1&2\end{pmatrix}.
\]
Thus
\[
        I_{C_3}\cong A_2.
\]

Let
\[
        \omega=e^{2\pi i/3},
        \qquad
        \OO=\ZZ[\omega].
\]
Evaluation at $\omega$ gives
\[
        \ZZ[C_3]/(1+g+g^2)\cong \ZZ[\omega].
\]
The norm is
\[
        N(m+n\omega)=m^2-mn+n^2.
\]
For the $A_2$ normalization
\[
        Q_{A_2}(m,n)=2(m^2-mn+n^2),
\]
the Epstein zeta is
\[
        Z_{A_2}(s)=\sum_{(m,n)\ne(0,0)}\left(2(m^2-mn+n^2)\right)^{-s}.
\]
Since $\ZZ[\omega]$ has six units and class number one,
\[
        Z_{A_2}(s)=2^{-s}6\zeta_{\QQ(\sqrt{-3})}(s)
        =2^{-s}6\zeta(s)L(s,\chi_{-3}).
\]

The Eisenstein--Leech note proves the adjacent lattice picture: the local Eisenstein plane is similar to $A_2$, twelve such planes assemble to the Niemeier lattice $N(12A_2)$, the Leech lattice is obtained as a Kneser/holy 3-neighbor, and the Leech theta series satisfies
\[
        \Theta_\Lambda=E_4^3-720\Delta,
        \qquad
        j=\Theta_\Lambda/\Delta+720
\]
\cite{EisensteinLeechNote}.  Thus the $A_2$ augmentation kernel is not isolated; it is the local plane of a well-known $24$-dimensional moonshine-adjacent lattice construction.

\section{CUE, Weyl denominators, and the derivative-moment Lie functor}

Let
\[
        \Phi_N(z)=\det(zI-U),
        \qquad U\sim\text{Haar}(U(N)).
\]
The compact Mellin note studies
\[
        M_N(k)=\mathbb E_{U(N)}|\Phi'_N(1)|^{2k},
\]
realizes it for fixed $N$ as a compact Mellin transform on $[0,N^24^{N-1}]$, and develops the $N=3$ phase-cover descent structure \cite{CompactMellinPeeling}.  Grover--Mezzadri--Simm study the more general derivative moments
\[
        M_{\mu,\nu}(z,N)=\mathbb E\left[
        \prod_{i=1}^{\ell(\mu)}\Lambda_N^{(\mu_i)}(z)
        \prod_{j=1}^{\ell(\nu)}\overline{\Lambda_N^{(\nu_j)}(z)}
        \right],
\]
with inner-disc exponent
\[
        \ell(\mu)\ell(\nu)+|\mu|+|\nu|
\]
and microscopic exponent
\[
        |\mu|+|\nu|+s^2
\]
when $\ell(\mu)=\ell(\nu)=s$ \cite{GroverMezzadriSimm}.

\subsection{Matrix-size $A_{N-1}$}

The Weyl integration formula for $U(N)$ involves
\[
        \prod_{1\le i<j\le N}(1-z_i/z_j)(1-z_j/z_i).
\]
The exponents are
\[
        e_i-e_j,
        \qquad i\ne j,
\]
the roots of $A_{N-1}$.  Thus fixed matrix size $N$ supplies an $A_{N-1}$ compact-torus root integral.

For $N=3$,
\[
        \Phi_3'(1)=3-2e_1+e_2,
\]
where
\[
        e_1=z_1+z_2+z_3,
        \qquad
        e_2=z_1z_2+z_1z_3+z_2z_3.
\]
The exact first moment is
\[
        M_3(1)=|3|^2+|-2|^2+|1|^2=14.
\]
The compact Mellin note also gives the exact constant-term checks
\[
        M_3(2)=375,
        \quad
        M_3(3)=15510,
        \quad
        M_3(4)=847280,
        \quad
        M_3(5)=55331424
\]
\cite{CompactMellinPeeling}.

\begin{theorem}[exact first derivative moment]
For every $N\ge1$,
\[
        M_N(1)=\mathbb E_{U(N)}|\Phi_N'(1)|^2
        =\sum_{m=1}^Nm^2
        =\frac{N(N+1)(2N+1)}6.
\]
\end{theorem}

\begin{proof}
Write
\[
        \Phi_N(z)=\sum_{j=0}^N(-1)^j e_j(U)z^{N-j}.
\]
Then
\[
        \Phi_N'(1)=\sum_{j=0}^{N-1}(-1)^j(N-j)e_j(U).
\]
The functions $e_j$ are the characters of the exterior powers $\wedge^j\CC^N$ and are orthonormal under Haar measure.  Hence
\[
        M_N(1)=\sum_{j=0}^{N-1}(N-j)^2=\sum_{m=1}^Nm^2.
\]
The final formula is the usual sum of squares.\end{proof}

\subsection{Moment-order $\mathfrak{sl}_{s+1}$}

For first derivative $s$-fold moments,
\[
        \mu=\nu=(1^s).
\]
Then
\[
        |\mu|=|\nu|=s,
\]
so the Grover--Mezzadri--Simm microscopic exponent is
\[
        s^2+2s=(s+1)^2-1=\dim\mathfrak{sl}_{s+1}.
\]
This can be made functorial.

Let
\[
        W=\CC^s.
\]
Define
\[
        \mathfrak g_s=\End(W)\oplus W\oplus W^\vee.
\]
Map
\[
        (A,u,\varphi)
        \longmapsto
        \begin{pmatrix}
        A&u\\
        \varphi&-\tr(A)
        \end{pmatrix}
        \in\mathfrak{sl}(W\oplus\CC).
\]

\begin{theorem}[CUE exponent Lie functor]
The above map is a vector-space isomorphism
\[
        \End(W)\oplus W\oplus W^\vee
        \cong
        \mathfrak{sl}_{s+1}(\CC).
\]
Under this isomorphism, the three CUE exponent sources
\[
        \ell(\mu)\ell(\nu)=s^2,
        \qquad
        |\mu|=s,
        \qquad
        |\nu|=s
\]
correspond to
\[
        \End(W),
        \qquad
        W,
        \qquad
        W^\vee.
\]
\end{theorem}

\begin{proof}
Every element of $\mathfrak{sl}(W\oplus\CC)$ has a unique block form
\[
        \begin{pmatrix}
        B&u\\
        \varphi&\lambda
        \end{pmatrix}
\]
with
\[
        \tr(B)+\lambda=0.
\]
Thus $\lambda=-\tr(B)$, and the element is uniquely determined by $(B,u,\varphi)$.  This proves the vector-space isomorphism.  The bracket is induced by the matrix commutator.\end{proof}

Thus:
\[
        s=1\Rightarrow \dim=3=\dim\mathfrak{sl}_2,
\]
\[
        s=2\Rightarrow \dim=8=\dim\mathfrak{sl}_3,
\]
\[
        s=3\Rightarrow \dim=15=\dim\mathfrak{sl}_4.
\]
The $A_2$ layer appears at fixed matrix size $N=3$ and at first-derivative moment order $s=2$.

\subsection{The phase-cover/SOS obstruction}

The compact Mellin note proves that in the first nontrivial finite layer $N=3$, integer kernels are Hermitian-square sums on a phase cover, while their descended quotient polynomials in variables
\[
        u=|\lambda|^2,
        \qquad
        v=|c|^2
\]
are governed by a quadratic trace recurrence and need not be real sums of squares.  In particular, the descended kernel $\Psi_2^\#(u,v)$ has a Gram matrix with negative principal minor, and hence is not an SOS over $\RR[u,v]$ \cite{CompactMellinPeeling}.

This is structurally the same quotient/lift pattern as split-zero support: the cover retains primitive decomposition data; the quotient forgets it; positivity or exact decomposition may fail to descend.  The precise common invariant is not that the quotient polynomial equals an $A_2$ form, but that both are governed by loss of a primitive lift whose virtual dimension-zero kernel is rank two in the $r=3$ case.

\section{Moonshine, Niemeier lattices, and the $Z_3$ triple}

Eguchi--Ooguri--Tachikawa observed that coefficients in the $K3$ elliptic genus decompose into dimensions of $M_{24}$ representations \cite{EguchiOoguriTachikawa}.  The umbral moonshine framework extends this to Niemeier lattices: the $23$ non-Leech Niemeier lattices, with full-rank ADE root systems of common Coxeter number, determine umbral groups and vector-valued mock modular forms; the pure $A$-type cases include root systems built from copies of $A_{\ell-1}$ \cite{ChengDuncanHarvey}.  Duncan's code construction realizes several umbral groups through distinguished linear codes over quotients of rings of integers and situates them inside the Leech/Conway world \cite{DuncanCodes}.

The $Z_3$-orbifold VOA theorem of Miyamoto gives a precise $C_3$ simple-current line.  Under the stated hypotheses, $V^G$ is semirigid, simple $V^G$-modules occur inside twisted or ordinary modules, and for solvable $G$ complete reducibility follows.  For lattice VOAs and order-three automorphisms, $V_L^\sigma$ is $C_2$-cofinite; the construction produces a moonshine VOA example and a $c=24$ VOA in Schellekens' list \cite{MiyamotoZ3}.

In the order-three construction, one has simple currents
\[
        W_0,
        \quad W_3,
        \quad W_6
\]
with
\[
        W_3\boxtimes W_3=W_6,
        \qquad
        W_3\boxtimes W_6=W_0.
\]
Thus
\[
        \ZZ\{W_0,W_3,W_6\}\cong\ZZ[C_3].
\]
The defect is
\[
        -\epsilon(aW_0+bW_3+cW_6)=-(a+b+c),
\]
so
\[
        -\epsilon(W_0+W_3+W_6)=-3,
\]
and
\[
        \ker\epsilon\cong A_2.
\]
This is the VOA realization of the triplet defect.

The Eisenstein--Leech note supplies the geometric lattice companion: the Eisenstein plane is the local $A_2$ geometry, twelve copies assemble to $N(12A_2)$, and the Leech lattice is obtained by a 3-neighbor/holy construction; the local order-three stabilizer forces cubic descent at the hexagonal CM point \cite{EisensteinLeechNote}.  This supports the structural chain
\[
        C_3\text{ augmentation}
        \Rightarrow A_2
        \Rightarrow \ZZ[\omega]
        \Rightarrow 12A_2
        \Rightarrow \Lambda_{\mathrm{Leech}}
        \Rightarrow \text{moonshine-adjacent }J\text{-function data}.
\]

\section{F-theory, $U(1)$ rank, and BSD as rank formalism}

F-theory geometrizes the type-IIB axio-dilaton
\[
        \tau_{IIB}=C_0+ie^{-\phi}
\]
as the complex modulus of an elliptic curve.  Vafa's $K3$ example has discriminant of degree $24$, corresponding to $24$ seven-brane degeneration points; the naive $U(1)^{24}$ count is reduced by the actual moduli constraints to $U(1)^{18}$ plus graviphoton contributions \cite{VafaFTheory}.  The same paper gives the local enhancement rule
\[
        A_{n-1}\text{ singularity}\quad\leadsto\quad SU(n)\text{ gauge symmetry},
\]
including examples with $A_4\mapsto SU(5)$ and $A_2\mapsto SU(3)$.

The defect construction aligns with this exactly:
\[
        r\text{ primitive }U(1)\text{ charges}
        \quad\leadsto\quad
        \ZZ^r\xrightarrow{-\Sigma}\ZZ,
\]
\[
        \ker(-\Sigma)=A_{r-1}
        \quad\leadsto\quad
        SU(r)\text{ root enhancement}.
\]
The repeated primitive $-1$ is the repeated $U(1)$ charge before root differences are imposed.  The roots are differences between such primitive charges.  This is the mathematical form of the stack-of-branes intuition: individual branes carry abelian gauge factors, while coincident or singular configurations produce nonabelian enhancement in the relative lattice.

Mordell--Weil rank gives the abelian side of elliptic-fibration gauge theory.  For an elliptic fibration with section,
\[
        \rk \MW(\pi)
\]
controls the number of independent abelian $U(1)$ factors, while reducible fibers contribute ADE root lattices.  BSD, on the arithmetic side, relates the Mordell--Weil rank of an elliptic curve over a number field to the analytic order of vanishing of its $L$-function; over $\QQ$, the rank $0$ and $1$ cases are theorem-level in the Gross--Zagier--Kolyvagin framework summarized in arithmetic-geometry sources \cite{ArithmeticGeometry}.

Thus the safe rank bridge is
\[
        \text{BSD analytic rank}
        \leftrightarrow
        \text{Mordell--Weil rank}
        \leftrightarrow
        U(1)\text{ gauge rank in elliptic fibrations},
\]
while the nonabelian root bridge is
\[
        A_{r-1}\text{ singularity/root lattice}
        \leftrightarrow
        SU(r)\text{ enhancement}.
\]
These are parallel structures through elliptic geometry.  They become a single theorem only after specifying an arithmetic elliptic fibration and its physical compactification data.

\section{Rank-one perturbations and the primitive $-1$ operator}

The primitive $-1$ operation also appears operator-theoretically as a rank-one boundary/evaluation perturbation.  The Connes--Consani--Moscovici zeta spectral triples paper constructs self-adjoint operators
\[
        D_{\log}^{(\lambda,N)}
        =D_{\log}^{(\lambda)}-|D_{\log}^{(\lambda)}\xi\rangle\langle\delta_N|
\]
as rank-one perturbations of a scaling spectral triple, and computes the regularized determinant
\[
        \det_{\mathrm{reg}}(D_{\log}^{(\lambda,N)}-z)
        =-i\lambda^{-iz}\widehat\xi(z),
\]
with zeros of $\widehat\xi$ lying on the real line under their hypotheses \cite{CCMSZ}.  The strategic convergence to the Riemann $\Xi$ function remains the missing step in that program; the relevant point here is structural: a rank-one boundary/evaluation correction is the operator-theoretic analogue of a single primitive defect sector.

This does not identify the split-zero construction with the spectral-triple construction.  It identifies the shared primitive operation: one boundary/evaluation vector contributes one rank-one correction, hence one primitive $-1$-type sector under the negative dimension/rank dictionary.

\section{Nielsen geometry and Epstein-complexity metrics}

Nielsen--Dowling geometry places a right-invariant Riemannian metric on a unitary group, decomposing Hamiltonians into easy and hard subspaces, typically one- and two-body terms versus higher-body terms.  The standard metric takes the form
\[
        \langle H,J\rangle
        =\frac{\tr(H_PJ_P)+q\tr(H_QJ_Q)}{2^n},
\]
and the geodesic equation determines optimal Hamiltonian evolution in this geometry \cite{NielsenDowling}.

The Weyl-operator construction supplies a finite Heisenberg basis.  For $r=3$, the nonidentity operators $W_{a,b}$ span $\mathfrak{sl}_3$.  If one lifts the labels $(a,b)\in\ZZ_3^2$ to the Eisenstein lattice $A_2$ and assigns penalty weights by
\[
        \lambda_{m,n}=Q_{A_2}(m,n)=2(m^2-mn+n^2),
\]
then the spectral zeta of the penalty spectrum is the $A_2$ Epstein zeta
\[
        Z_{A_2}(s)=2^{-s}6\zeta(s)L(s,\chi_{-3}).
\]
This gives a concrete construction of an Epstein-complexity metric: use a right-invariant Nielsen metric diagonalized by Weyl operators, with penalties assigned by the lifted root-lattice quadratic form.

For general $r$, the analogous construction uses the $A_{r-1}$ root lattice and the Epstein zeta $Z_{A_{r-1}}(s)$.  The $r=3$ case is special because $A_2$ is the Eisenstein lattice and therefore admits the CM/Dedekind factorization above.

\section{Consolidated theorem}

\begin{theorem}[unified primitive-defect structure]
The following constructions realize the same primitive defect mechanism.
\begin{enumerate}[label=(\roman*)]
\item Split-zero support: Boolean support in $G(R)^I$ linearizes to
\[
        \ZZ^I\xrightarrow{-\Sigma}\ZZ.
\]
\item Root lattice: for $|I|=r$,
\[
        \ker(-\Sigma)=A_{r-1}.
\]
\item Primitive $U(1)$ charges: $\ZZ^r$ is the character lattice of $U(1)^r$, and $A_{r-1}$ is the difference lattice of relative charges.
\item Semilocal boundary theory: each finite-prime boundary unit and each archimedean minimal stabilizer projection contributes $-1$ under the connecting map; in the $S=\{p,\infty\}$ visible boundary, the total triplet class has defect $-3$ and kernel $A_2$.
\item Fusion: for a cyclic simple-current line,
\[
        \ZZ[C_r]\xrightarrow{-\epsilon}\ZZ
\]
has kernel $A_{r-1}$.
\item Modular diagonalization: in the pointed cyclic fusion algebra, $-\epsilon$ is the negative trivial Fourier/$S$-row functional.
\item Discrete-torsion double: $C_r\times\widehat C_r$ with the Heisenberg cocycle yields $M_r(\CC)$; the nonidentity Weyl operators span $\mathfrak{sl}_r$.
\item CUE: fixed matrix size $N$ has $A_{N-1}$ Weyl denominator, and first-derivative moment order $s$ has exponent $s^2+2s=\dim\mathfrak{sl}_{s+1}$ realized by
\[
        \End(\CC^s)\oplus\CC^s\oplus(\CC^s)^\vee.
\]
\item Eisenstein/Leech/moonshine: the $r=3$ cyclic augmentation ideal is $A_2$, identified with the Eisenstein plane; twelve copies produce the $12A_2$ Niemeier context adjacent to Leech and $Z_3$-orbifold moonshine.
\item F-theory: $r$ primitive abelian factors form the $U(1)^r$ side, while the difference lattice $A_{r-1}$ gives the local $SU(r)$ enhancement root system.
\end{enumerate}
\end{theorem}

\begin{proof}
Items (i)--(iii) follow from Theorem 1.1 and the split-zero support map.  Item (iv) follows from the finite-prime boundary computation and the archimedean fixed-point computation.  Item (v) is the augmentation ideal calculation for $\ZZ[C_r]$.  Item (vi) is Fourier diagonalization of the pointed fusion matrices.  Item (vii) is the Weyl-operator theorem.  Item (viii) combines the Weyl denominator calculation with the Grover--Mezzadri--Simm exponent and the block isomorphism with $\mathfrak{sl}_{s+1}$.  Item (ix) is the $C_3$ augmentation/Eisenstein identification plus the cited Niemeier/Leech and Miyamoto $Z_3$-orbifold constructions.  Item (x) is the standard $A_{r-1}\mapsto SU(r)$ singularity enhancement and Mordell--Weil rank dictionary in elliptic-fibration geometry.\end{proof}

\section{Further problems}

The theorem suggests several precise next constructions.

\begin{enumerate}[label=(\arabic*)]
\item Construct an explicit functor from the CUE $N=3$ phase-cover trace algebra to the $A_2$ defect lattice, not merely a quotient/lift analogy.  The expected target is a category whose objects are phase-cover Hermitian-square decompositions and whose dimension-zero kernel is the $A_2$ virtual support lattice.

\item Construct the full split-zero matrix semiring $G(M_r(\CC))$ and its support-diagram category for Weyl-operator expansions.  The basic theorem must separate supported zero coefficients, unsupported absent sectors, and trace removal.

\item Extend the semilocal finite-prime edge calculation from the one-zero visible truncation to arbitrary iterated boundary strata.  The predicted primitive map is again negative rank on each boundary circle, but higher-codimension strata should produce higher cubical incidence maps.

\item Build the $A_2$ Eisenstein Epstein-complexity metric on $SU(3)$ explicitly in Nielsen geometry, compute its geodesic equation in Weyl coordinates, and compare its spectral zeta with $Z_{A_2}$.

\item Develop the $A_2\oplus A_1\oplus A_0$ Standard-Model root skeleton inside the support-defect formalism, then impose the finite central quotient and hypercharge lattice as additional data.

\item Specify an arithmetic elliptic fibration for which the Mordell--Weil rank, BSD analytic rank, and F-theory $U(1)$ rank can be compared as a single object rather than parallel rank formalisms.
\end{enumerate}

\begin{thebibliography}{99}

\bibitem{SplitZeroNotes}
\textit{Split-zero globalization and secondary notes}.  Unpublished project notes, 2026.

\bibitem{ConnesConsaniKnots}
A. Connes and C. Consani, \textit{Knots, primes and the adele class space}; and related semilocal computations in \textit{Knots, primes and class field theory}.

\bibitem{DVVV}
R. Dijkgraaf, C. Vafa, E. Verlinde, and H. Verlinde, \textit{The operator algebra of orbifold models}, Commun. Math. Phys. 123 (1989), 485--526.

\bibitem{DixonHarveyVafaWitten}
L. Dixon, J. Harvey, C. Vafa, and E. Witten, \textit{Strings on orbifolds}, Nucl. Phys. B261 (1985), 678--686.

\bibitem{MiyamotoZ3}
M. Miyamoto, \textit{A $\mathbb Z_3$-orbifold theory of lattice vertex operator algebra and $\mathbb Z_3$-orbifold constructions}, arXiv:1003.0237.

\bibitem{EguchiOoguriTachikawa}
T. Eguchi, H. Ooguri, and Y. Tachikawa, \textit{Notes on the K3 surface and the Mathieu group $M_{24}$}, arXiv:1004.0956.

\bibitem{DuncanCodes}
J. F. R. Duncan, \textit{Umbral moonshine and linear codes}, 2013.

\bibitem{ChengDuncanHarvey}
M. C. N. Cheng, J. F. R. Duncan, and J. A. Harvey, \textit{Umbral moonshine and the Niemeier lattices}, arXiv:1307.5793.

\bibitem{EisensteinLeechNote}
\textit{Eisenstein planes, the Leech lattice, the cubic mirror coordinate, genus-two theta coefficients, and the Euler-hypergeometric shadow}.  Unpublished note, 2026.

\bibitem{GroverMezzadriSimm}
A. Grover, F. Mezzadri, and N. Simm, \textit{Higher order derivative moments of CUE characteristic polynomials and the Riemann zeta function}, arXiv:2604.03051.

\bibitem{CompactMellinPeeling}
\textit{Compact Mellin peeling, transform algebras, and the first noninteger layers of CUE derivative moments}.  Unpublished note, 2026.

\bibitem{VafaFTheory}
C. Vafa, \textit{Evidence for F-theory}, arXiv:hep-th/9602022.

\bibitem{StromingerVafa}
A. Strominger and C. Vafa, \textit{Microscopic origin of the Bekenstein-Hawking entropy}, arXiv:hep-th/9601029.

\bibitem{WGC}
N. Arkani-Hamed, L. Motl, A. Nicolis, and C. Vafa, \textit{The string landscape, black holes and gravity as the weakest force}, arXiv:hep-th/0601001.

\bibitem{OoguriVafa}
H. Ooguri and C. Vafa, \textit{On the geometry of the string landscape and the swampland}, arXiv:hep-th/0605264.

\bibitem{ArithmeticGeometry}
H. Darmon et al., \textit{Arithmetic Geometry}, Clay Mathematics Proceedings, vol. 8, AMS/CMI, 2009.

\bibitem{NielsenDowling}
M. R. Dowling and M. A. Nielsen, \textit{The geometry of quantum computation}, arXiv:quant-ph/0701004.

\bibitem{CCMSZ}
A. Connes, C. Consani, and H. Moscovici, \textit{Zeta spectral triples}, arXiv:2511.22755.

\end{thebibliography}

\end{document}
